\documentclass[10pt,a4paper]{amsart}
\usepackage[margin=19mm]{geometry}
\usepackage{amsmath,amssymb,mathtools}
\usepackage[scr=rsfs,cal=euler]{mathalfa}
\usepackage{xfrac}
\usepackage[unicode,hidelinks,bookmarks=false]{hyperref}
\newcommand{\ie}{i.e.\ }
\newcommand{\cf}{cf.\ }
\newcommand{\resp}{respectively\ }
\newcommand{\Subsection}{\subsection}
\providecommand{\nobreakdash}{\nobreak\hbox{-}}
\setlength{\emergencystretch}{2em}
\multlinegap0pt
\usepackage{bm}
\usepackage{bbm}
\usepackage{dsfont}
\usepackage{xspace}
\usepackage{cancel}
\usepackage{mathtools}
\usepackage{amsthm}
\makeatletter
\@ifundefined{theorem}{\newtheorem{theorem}{Theorem}}{}
\@ifundefined{lemma}{\newtheorem{lemma}[theorem]{Lemma}}{}
\@ifundefined{proposition}{\newtheorem{proposition}[theorem]{Proposition}}{}
\makeatother
\title[The Minimum Subgraph Complementation Problem]{The Minimum Subgraph Complementation Problem}
\author{Juan Guti\'errez, Sagartanu Pal}
\date{Source: arXiv:2512.23687v2 (CC BY 4.0)}
\begin{document}
\maketitle

\noindent\emph{Source and licence.} The Minimum Subgraph Complementation Problem. Version arXiv:2512.23687v2 (CC BY 4.0), distributed under the Creative Commons Attribution 4.0 International licence, as recorded in the arXiv licence metadata for this exact version and shown on the abstract page https://arxiv.org/abs/2512.23687v2. Full rights notice: licenses/arxiv-2512.23687-CC-BY-4.0.txt.\par\smallskip

\noindent\textit{Editorial scope.} Counted unit: the Lemma whose source label is \texttt{lemma:forest-to-degeneracy} in the article (source0.tex lines 1119--1126, source label \texttt{lemma:forest-to-degeneracy}), delivered together with the article's own complete original proof of it (source0.tex lines 1127--1186, one \texttt{proof} environment of 60 lines). No mathematics is written, repaired or re-derived: statement and proof are the article's own text line for line. Required context retained uncounted: prop:GinmathcalD\_kthenaandb (lines 1084--1090); prop:Sgeqk (lines 1095--1098). Every definition, display and label of those lines is retained unchanged. Omitted from this package and not redistributed: 1--1083, 1091--1094, 1099--1118, 1187--1757. Nothing inside a retained excerpt is omitted or altered: every retained body is delivered exactly as the article's lines are written, so no omission receipt of any kind accompanies this package. Citations: the retained excerpts carry no citation command at all, so no bibliography entry of the article is transcribed and no reference is redistributed. Reference adaptation: none was needed and none was made; every reference inside the delivered unit resolves inside it, so no unresolved reference or citation is produced. Printed slips kept literal: no printed word, symbol, hypothesis or number of the counted statement or of its proof was altered. Layout and class: the article's own class is \texttt{article} with packages including fontenc, inputenc, microtype, fullpage, comment, amsmath, amssymb, amsthm, graphicx, subcaption, booktabs, tikz, algorithm, algpseudocode; the wrapper is the lane's pinned \texttt{amsart} editorial wrapper at 10pt on A4, which restates the theorem-like environments the retained text needs and adds the article's own macro definitions that the retained text needs (none), with extra wrapper packages (bm, bbm, dsfont, xspace, cancel, mathtools). The delivered numbering is the unit's own. Assets: no retained span contains a figure environment, a graphics-inclusion command, a PDF-page inclusion, a table or a TikZ picture; no image file is copied into this package, the package declares no asset dependency and no image file is redistributed with it. Licence: version arXiv:2512.23687v2 of this article is distributed under the Creative Commons Attribution 4.0 International licence, as recorded in the arXiv licence metadata for this exact version and shown on the abstract page https://arxiv.org/abs/2512.23687v2. A full rights notice is delivered at licenses/arxiv-2512.23687-CC-BY-4.0.txt. Characters: every retained excerpt is pure ASCII, so the delivered engine, XeLaTeX with its default Latin Modern text font, reports no missing character.
\begin{proposition} \label{prop:GinmathcalD_kthenaandb}
A graph $G$ belongs to $\mathcal{D}_k$ if and only if
\begin{itemize}
    \item[(a)] every induced subgraph of $G$ has a vertex of degree at most $k$, and
    \item[(b)] there exists an induced subgraph of $G$ in which every vertex has degree at least $k$.
\end{itemize}
\end{proposition}

\begin{proposition} \label{prop:Sgeqk}
Let $k \geq 2$.
Let $G$ be a forest. If $S$ is a solution, then $|S| \geq k$.
\end{proposition}

\begin{lemma}
\label{lemma:forest-to-degeneracy}
Every graph in $\mathcal{F}_k$ is complementable to $\mathcal{D}_k$ for any fixed $k\geq 2$.
 Let $G \in \mathcal{F}_k$
and $S$ be a minimum solution.
    If $G$ has a sibling set of size $k$,
    then $|S|=k$. Otherwise, $|S|=k+1$.
\end{lemma}

\begin{proof}
Let $G \in \mathcal{F}_k$.
%i had to modify because we need to allow empty sets as solutions
Suppose for a moment that $G$ has a sibling set of size $k$, say $S$.
We will show that $G'=G \oplus S$ has degeneracy $k$. As every vertex in $G'[S \cup \{u\}]$, where $u$
 is the common neighbor of the sibling set,
has degree $k$, part $(b)$ of Proposition \ref{prop:GinmathcalD_kthenaandb} follows.
Let $X \subseteq V(G')$. We need to show that there exists a vertex in $G'[X]$ with degree at most $k$.
If $X \subseteq S$, then any vertex in $G'[X]$ has degree at most $k$, so part $(a)$ of Proposition \ref{prop:GinmathcalD_kthenaandb} holds.
Otherwise, $X$ has a vertex in $\overline{S}$.

Note that $G'[X \cap \overline{S}]$ is a forest.
If $|X \cap \overline{S}|=1$, then the only vertex in
$X \cap \overline{S}$ has at most $k$ neighbors in $G[X \cap S]$, and therefore also in $G'[X]$ and part (a) of Proposition \ref{prop:GinmathcalD_kthenaandb} follows.  
Hence, we may assume that $|X \cap \overline{S}|>1$, which implies
that $G'[X \cap \overline{S}]$ has two vertices, say $v$ and $w$, of degree at most one.
Suppose by contradiction that none of them has degree at most $k$ in $G'[X]$. In that case, both vertices have degree at least $k+1$. As both vertices have degree at most one in $G'[X \cap \overline{S}]$, it implies that both $v$ and $w$ are adjacent, in $G'$, to all vertices in $S$.
Note that this implies that, in $G$, they are also adjacent to all vertices in $S$, a contradiction to the fact that $G$ is a forest since $k\geq 2$.
Now, by Proposition \ref{prop:Sgeqk}, $S$ is a minimum solution and the proof follows.


Finally, suppose that $G$ does not have a sibling set of size $k$. As $G$ has at least $2k+2$ vertices, it has an independent set, say $S$ of size $k+1$.
We will show that $G'=G \oplus S$ is a graph with degeneracy $k$. As every vertex in $G'[S]$ has degree $k$, then part $(b)$ of Proposition \ref{prop:GinmathcalD_kthenaandb} follows.
Let $X \subseteq V(G')$. We need to show that there exists a vertex in $G'[X]$ with degree at most $k$.
If $X \subseteq S$, then any vertex in $G'[X]$ has degree at most $k$, so part $(a)$ of Proposition \ref{prop:GinmathcalD_kthenaandb} holds.
Otherwise, $X$ has a vertex in $\overline{S}$.
Let $v$ be a vertex of minimum degree in $G'[X \cap \overline{S}]$. As $G'[X \cap \overline{S}]$ is a forest, such a vertex has degree at most one in
$G'[X \cap \overline{S}]$.
As there is no subset of $k$ siblings in $G$,
$v$ is adjacent to at most $k-1$ vertices in $S$.
Hence, $v$ has at most $k$ neighbors in $G'[X]$ and part $(a)$ of Proposition \ref{prop:GinmathcalD_kthenaandb} follows.
%Therefore, every graph in $\mathcal{F}_k$ is complementable to $\mathcal{D}_k$ for any $k\geq 2$. 
%Next, we will show that $S$ has minimum size over all subsets.
Next, we will show that $S$ is a minimum solution.

Suppose first that $k=2$. Then $G$ has at least 6 vertices.
 As $G$ does not have a sibling set of size 2, and $G$ has no cycles, 
each component of $G$ is either a $K_2$ or a $K_1$.
Then for any subset $S$ of size 2, in $G'=G\oplus S$,
Proposition \ref{prop:GinmathcalD_kthenaandb}(b) is not true.
Therefore, the size of the minimum solution is 3, i.e., $k+1$.


Consider now the case $k\geq 3$.
Suppose by contradiction that there exists $R \subseteq V(G)$ such that $G'=G \oplus R \in \mathcal{D}_k$ and $|R| < |S| = k+1$.
By Proposition \ref{prop:Sgeqk}, $|R|=k$.
By Proposition \ref{prop:GinmathcalD_kthenaandb}(b), there exists
    $X \subseteq V(G')$  such that every vertex in $G'[X]$ has degree at least $k$, which implies 
    that $X \setminus R \neq \emptyset$.
    We will show that $G' [X \cap \overline{R}]$ has no isolated vertices.
For a contradiction, suppose $G' [X \cap \overline{R}]$ has an isolated vertex $v$.
Since $|R|=k$ and any vertex in $G'[X]$ has degree at least $k$, $v$ is 
adjacent to all the vertices of $R$ in $G'[X]$, which implies that $R$ is a sibling set in $G$, a contradiction.
Note that 
$G' [X \cap \overline{R}]$ is a forest.
Let $u$ and $v$ be two vertices of degree 1 in the same component of $G' [X \cap \overline{R}]$.
Then both $u$ and $v$ are adjacent to at least $k-1$ vertices of
$R$ in $G'[X]$. Since $k\geq 3$, $u$ and $v$ have a common neighbor in $X\cap R$, 
which implies there exists a cycle in $G$, a contradiction.
\end{proof}

\end{document}
